\subsection*{小结}

\begin{table}[htbp]\centering\small
\begin{tabularx}{\linewidth}{@{}lX@{}}
\toprule
| \textbf{主题} & \textbf{要点} | \\
\midrule
| --- & --- | \\
| 二维张量收缩 & 在分块状操作数上做 C += A × B —— 通用的内层核函数 | \\
| 四步抽象 & \textbf{填充 → 加载 → 相乘 → 存储}；编译器为各目标平台处理片元布局 | \\
| 扩大协作规模 & 语法保持不变 | \\
| 布局约定 & \co@mma.row.row@、\co@mma.row.col@ 等 | \\
| 稀疏与量化 & \co@.sp@ 增加元数据操作数；\co@.scale@ 融合逐分块反量化 | \\
| 重排与流水线 & \co@mma.load.swiz<N>@ 匹配带重排的共享布局 | \\
\bottomrule
\end{tabularx}
\end{table}

收缩执行得更快，但装载和计算仍然轮流进行——3.8 节会引入\textbf{角色专门化}，让不同的线程组各司其职：一组装载数据，另一组做计算，使内存访问和数学运算在时间上重叠。
